\section{Indexing in MongoDB}
\label{sec:indexing}

\subsection{Theory (in class)}
\label{sec:index-theory}

This subsection summarizes the file-organization and indexing theory of the course
(Chapter~2, based on \cite{elmasri2016,dbms-ch2}). It starts from how records are stored on disk,
reviews the basic file organizations, and then builds up from single-level indexes to multilevel
indexes and B$^+$-trees. It is the reference point for the comparison with MongoDB in
\Cref{sec:index-compare}.

\subsubsection{Records, blocks and the cost model}
A \emph{file} is a sequence of \emph{records}, and each record is a collection of field values.
Records may have fixed or variable length. Records are stored in disk \emph{blocks}, and the
block is the unit of transfer between disk and memory. Storing several records in one block is
called \emph{blocking}, and the \emph{blocking factor} is the number of records per block:
\[
bfr = \left\lfloor \frac{B}{R} \right\rfloor, \qquad
b = \left\lceil \frac{r}{bfr} \right\rceil ,
\]
where $B$ is the block size, $R$ the record size, $r$ the number of records and $b$ the number
of blocks of the file. Records may be \emph{unspanned} (a record never crosses a block boundary,
usual for fixed-length records) or \emph{spanned} (a record may continue in the next block, usual
for variable-length records). Because reading a block from disk is far more expensive than
processing it in memory, the cost of an operation is measured by the \textbf{number of block
accesses}. The typical file operations are \textsc{open}, \textsc{find}, \textsc{findnext},
\textsc{read}, \textsc{insert}, \textsc{delete}, \textsc{modify}, \textsc{close},
\textsc{reorganize} and \textsc{read\_ordered}.

\subsubsection{Basic file organizations}
\label{sec:file-org}
\begin{description}
  \item[Unordered (heap) file.] New records are appended at the end of the file, so insertion is
        very efficient. Searching requires a \emph{linear search}, which reads half of the blocks
        on average ($b/2$), and reading the records in the order of a field requires sorting.
  \item[Ordered (sequential) file.] Records are kept sorted by an \emph{ordering field}. A
        \emph{binary search} on the ordering field needs about $\lceil\log_2 b\rceil$ block
        accesses, and reading records in order is efficient. Insertion is expensive because the
        record must be placed in the right position; a separate unordered \emph{overflow} file is
        often used for new records and merged periodically into the main file.
  \item[Hashed file (external hashing).] The blocks are divided into $M$ buckets, and a record
        with hash-key value $K$ is stored in bucket $h(K)$. Search on the hash key is very
        efficient (about one block access). A \emph{collision} happens when a bucket is full;
        it is resolved by open addressing, chaining with an overflow area, or multiple hashing.
        A hash file is typically kept 70--80\% full and needs a hash function that spreads
        records uniformly. Its drawbacks are the fixed number of buckets $M$ when the file grows
        or shrinks, and the inefficiency of ordered access on the hash key.
\end{description}
\Cref{tab:file-org} summarizes the average cost of finding one record. Only the ordering
(or hash) field benefits from the organization: a search on any other field still scans the file.
Indexes were introduced to give efficient access paths on those other fields as well.

\begin{table}[H]
\centering
\small
\begin{tabularx}{\textwidth}{@{}l>{\raggedright\arraybackslash}X>{\raggedright\arraybackslash}X@{}}
\toprule
\textbf{Organization} & \textbf{Search on the organizing field} & \textbf{Insertion} \\
\midrule
Unordered (heap) & Linear: $\approx b/2$ block accesses & Cheap: append to the end \\
Ordered & Binary search: $\lceil\log_2 b\rceil$ & Expensive: keep the order (or use an overflow file) \\
Hashed & About 1 (without overflow) & Cheap, but overflow may occur \\
\bottomrule
\end{tabularx}
\caption{Average cost of accessing a record under the three file organizations.}
\label{tab:file-org}
\end{table}

\subsubsection{Indexes as access paths}
A \emph{single-level index} is an auxiliary file that makes it more efficient to find records in
the data file. It is usually specified on one field of the file (the \emph{indexing field}) and
is a file of entries $\langle K(i), P(i)\rangle$, \emph{ordered} by the field value $K(i)$, where
$P(i)$ is a pointer to a block or to a record. The index is called an \emph{access path} on that
field. Because an index entry holds only the key and a pointer, it is much smaller than a data
record, so the index file occupies considerably fewer blocks than the data file. A binary search
on the (ordered) index file therefore touches few blocks and yields a pointer to the record, after
which one more block access fetches the data.

Indexes are further characterized as \emph{dense} or \emph{sparse}. A \textbf{dense} index has an
entry for every search-key value (hence for every record) of the data file. A \textbf{sparse}
(non-dense) index has entries for only some of the search values. A sparse index is smaller, but it
requires the data file to be ordered on the indexing field so that the record can be located from
the nearest preceding entry.

\subsubsection{Types of single-level ordered indexes}
\label{sec:single-level}
\paragraph{Primary index.}
A primary index is defined on a data file that is \emph{ordered on a key field}. It has one entry
per \emph{block} of the data file, made of the key of the first record of the block (the
\emph{block anchor}) and a pointer to that block. The number of entries therefore equals the
number of data blocks, and the index is sparse. To find a record with key $K$, binary-search the
index for the last entry whose anchor is $\le K$, then read that data block and search inside it.

\paragraph{Clustering index.}
A clustering index is defined on a data file that is \emph{ordered on a non-key field}, so that
records with the same value are stored together (clustered). It has one entry per \emph{distinct
value} of the field, pointing to the first block that contains records with that value. The index
is sparse. A file has at most one physical ordering, hence it can have at most one primary index
\emph{or} one clustering index, but not both.

\paragraph{Secondary index.}
A secondary index provides an additional access path on a field by which the data file is
\emph{not} ordered. The indexing field may be a \emph{key} (unique values) or a \emph{non-key}
(duplicate values), and a file can have many secondary indexes.
\begin{itemize}
  \item On a key field, the index is \emph{dense}: one entry per record, because records cannot
        be located from neighboring entries when the file is unordered on that field.
  \item On a non-key field, there are three options: (1) a duplicate entry for each record, (2)
        a variable-length list of pointers $\langle P(i,1),\dots,P(i,k)\rangle$ inside the entry,
        or (3) one entry per distinct value together with an extra level of indirection (blocks
        of record pointers). Option~3 is the most commonly used.
\end{itemize}

\Cref{tab:single-level} compares the three types.

\begin{table}[H]
\centering
\small
\begin{tabularx}{\textwidth}{@{}l>{\raggedright\arraybackslash}X>{\raggedright\arraybackslash}X>{\raggedright\arraybackslash}X@{}}
\toprule
 & \textbf{Primary} & \textbf{Clustering} & \textbf{Secondary} \\
\midrule
Data file ordered on indexing field & Yes & Yes & No \\
Indexing field is a key & Yes & No & Key or non-key \\
Number of index entries & One per data block & One per distinct value & One per record (key), or one per distinct value (non-key, option 3) \\
Dense / sparse & Sparse & Sparse & Dense (key); dense or sparse (non-key) \\
How many per file & At most one primary \emph{or} one clustering & (see left) & Many \\
\bottomrule
\end{tabularx}
\caption{Single-level ordered indexes.}
\label{tab:single-level}
\end{table}

\subsubsection{Cost of searching: worked examples}
\label{sec:index-examples}
Consider a file with $r = 300{,}000$ fixed-size, unspanned records of $R = 100$ bytes on a disk
with block size $B = 4{,}096$ bytes. Then
\[
bfr = \left\lfloor \frac{4096}{100} \right\rfloor = 40, \qquad
b = \left\lceil \frac{300000}{40} \right\rceil = 7{,}500 .
\]

\paragraph{Without an index.} On the ordered file, a binary search needs
$\lceil \log_2 7500 \rceil = 13$ block accesses. On an unordered file, a linear search needs on
average $b/2 = 3{,}750$ accesses.

\paragraph{Primary index.} Let the ordering key have $V = 9$ bytes and block pointers $P = 6$
bytes. An index entry has $R_i = 9 + 6 = 15$ bytes, so
$bfr_i = \lfloor 4096/15 \rfloor = 273$ entries per block. The index has $r_i = 7{,}500$ entries
(one per data block) and occupies $b_i = \lceil 7500/273 \rceil = 28$ blocks. A binary search on the
index costs $\lceil \log_2 28 \rceil = 5$ accesses, plus 1 access for the data block:
\textbf{6 block accesses} in total, against 13.

\paragraph{Clustering index.} Suppose the same file is ordered on \code{Zipcode} with 1{,}000
distinct values, and an entry has $5 + 6 = 11$ bytes. Then $bfr_i = \lfloor 4096/11 \rfloor = 372$,
$r_i = 1{,}000$ and $b_i = \lceil 1000/372 \rceil = 3$ blocks, so a binary search on the index needs
$\lceil \log_2 3 \rceil = 2$ accesses before reading the blocks that hold the records of that
zip code.

\paragraph{Secondary index on a key field.} For a non-ordering key field of the same file, the
index is dense: $r_i = 300{,}000$ entries, $b_i = \lceil 300000/273 \rceil = 1{,}099$ blocks. A
binary search costs $\lceil \log_2 1099 \rceil = 11$ accesses, plus 1 for the record:
\textbf{12 block accesses}, against $3{,}750$ for a linear search.

\subsubsection{Multilevel indexes}
\label{sec:multilevel}
Since a single-level index is itself an ordered file, we can create a primary index \emph{on the
index}: the original index is the \emph{first level} and the new index is the \emph{second level}.
The process is repeated, creating a third, fourth, \dots\ level, until all entries of the top level
fit in one block. A multilevel index can be built on any first-level index (primary, clustering or
secondary) provided it spans more than one block. The number of blocks at each level is
\[
b_1 = \left\lceil \frac{r_1}{bfr_i} \right\rceil, \quad
b_{j} = \left\lceil \frac{b_{j-1}}{bfr_i} \right\rceil, \quad
\text{stop at the top level } t \text{ where } b_t = 1,
\]
and a search reads one block per level plus the data block: $t + 1$ accesses. The factor $bfr_i$ is
the \emph{fan-out}, which is why large fan-out keeps the number of levels small. A two-level
primary index resembles the \emph{ISAM} (Indexed Sequential Access Method) organization.

For the secondary index of the example: $b_1 = 1{,}099$, $b_2 = \lceil 1099/273 \rceil = 5$ and
$b_3 = \lceil 5/273 \rceil = 1$, so $t = 3$ and a search costs $t + 1 = 4$ block accesses, instead
of 12 with binary search on the single-level index.

A multilevel index is a form of \emph{search tree}. Its weakness is maintenance: every level is
an ordered file, so inserting or deleting an entry may require shifting entries at several levels.
This is the ``severe problem'' that dynamic structures solve.

\subsubsection{B-trees and B$^+$-trees}
\label{sec:btree}
Most multilevel indexes use B-trees or B$^+$-trees because of the insertion and deletion problem.
They are search trees in which each node corresponds to a disk block and is kept between half full
and completely full, which leaves free space for new entries without reorganizing the whole index.

\paragraph{B-tree versus B$^+$-tree.}
In a B-tree, pointers to data records appear at \emph{all} levels of the tree. In a B$^+$-tree,
pointers to data records appear \emph{only in the leaf nodes}; internal nodes contain only search
values and tree pointers, and the leaf nodes are linked together in key order. Consequently an
internal B$^+$-tree node can hold more entries, so a B$^+$-tree can have fewer levels (a higher
capacity) than the corresponding B-tree, and ordered or range access is efficient by following the
leaf links. These properties explain why the B$^+$-tree is the standard index structure of
commercial DBMSs.

\begin{table}[H]
\centering
\small
\begin{tabularx}{\textwidth}{@{}l>{\raggedright\arraybackslash}X>{\raggedright\arraybackslash}X@{}}
\toprule
 & \textbf{B-tree} & \textbf{B$^+$-tree} \\
\midrule
Data pointers & In nodes at all levels & Only in leaf nodes \\
Internal nodes & Keys, data pointers and tree pointers & Keys and tree pointers only \\
Leaf nodes & Same structure as internal nodes & Key and data-pointer pairs, plus a pointer to the next leaf \\
Fan-out and height & Lower fan-out, possibly more levels & Higher fan-out, fewer levels \\
Range / ordered scan & Requires traversing the tree & Follow the linked leaves \\
\bottomrule
\end{tabularx}
\caption{Comparison of B-trees and B$^+$-trees.}
\label{tab:btree-vs-bplus}
\end{table}

\paragraph{Structure of a B$^+$-tree node.}
Let $p$ be the order of an \emph{internal} node: it has up to $p$ tree pointers and up to $p-1$
search values $K_1 < K_2 < \dots$, with the subtree at the $i$-th pointer holding values $\le K_i$
(and greater than $K_{i-1}$). A \emph{leaf} node holds up to $p_{leaf}$ pairs
$\langle K, P_r\rangle$ with record pointers, followed by a pointer to the next leaf. Every node
except the root is at least half full, and all leaves are at the same level, so the tree is always
\emph{balanced} and every search visits the same number of nodes.

\begin{figure}[H]
\centering
\begin{tikzpicture}[
  nd/.style={draw, rounded corners=1pt, minimum height=7mm, inner sep=5pt, font=\small\ttfamily},
  >=Stealth]
  \node[nd, fill=blue!8]  (root) at (0,0) {9 \quad 17};
  \node[nd, fill=green!8] (l1) at (-4,-2) {1 \quad 5 \quad 7};
  \node[nd, fill=green!8] (l2) at (0,-2)  {9 \quad 12 \quad 15};
  \node[nd, fill=green!8] (l3) at (4,-2)  {17 \quad 20 \quad 25};
  \draw[->] ([xshift=-4mm]root.south) -- (l1.north);
  \draw[->] (root.south) -- (l2.north);
  \draw[->] ([xshift=4mm]root.south) -- (l3.north);
  \draw[->, dashed] (l1.east) -- (l2.west);
  \draw[->, dashed] (l2.east) -- (l3.west);
  \foreach \n in {l1,l2,l3}
    \draw[->, thick] (\n.south) -- ++(0,-0.6) node[below, font=\scriptsize] {data records};
\end{tikzpicture}
\caption{A small B$^+$-tree: an internal node (blue) and linked leaf nodes (green, dashed arrows) with pointers to the data records.}
\label{fig:bplus}
\end{figure}

\paragraph{Order and capacity.}
Each node must fit in one block. With $B = 512$ bytes, key size $V = 9$, block pointer $P = 6$ and
record pointer $P_r = 7$, an internal node needs $p\cdot P + (p-1)\cdot V \le B$, i.e.\
$15p \le 521$, so $p = 34$. A leaf needs $p_{leaf}\,(P_r + V) + P \le B$, i.e.\ $16\,p_{leaf} \le 506$,
so $p_{leaf} = 31$. Assuming nodes are on average 69\% full, an internal node has about
$34 \times 0.69 \approx 23$ pointers (22 keys) and a leaf about $31 \times 0.69 \approx 21$ record
pointers. \Cref{tab:capacity} shows how the number of entries grows level by level.

\begin{table}[H]
\centering
\small
\begin{tabular}{@{}lrrr@{}}
\toprule
\textbf{Level} & \textbf{Nodes} & \textbf{Index entries} & \textbf{Pointers} \\
\midrule
Root & 1 & 22 & 23 \\
Level 1 & 23 & $23\times 22 = 506$ & $23^2 = 529$ \\
Level 2 & 529 & $529\times 22 = 11{,}638$ & $23^3 = 12{,}167$ \\
Leaf level & 12,167 & $12{,}167\times 21 = 255{,}507$ & -- \\
\bottomrule
\end{tabular}
\caption{Average capacity of a three-level B$^+$-tree (nodes 69\% full).}
\label{tab:capacity}
\end{table}

A three-level B$^+$-tree thus indexes about 255{,}507 records, compared with 65{,}535 entries for
the corresponding B-tree, and a search costs only the number of levels plus one data-block access.

\paragraph{Insertion.}
The new entry is inserted in the proper leaf. If the leaf has room, nothing else changes. If the
leaf is full it \emph{overflows} and is split: a new leaf is created, the first
$j = \lceil (p_{leaf}+1)/2 \rceil$ entries stay in the original leaf, the remaining entries move to
the new leaf, the $j$-th search value is \emph{replicated} in the parent, and the parent receives an
extra pointer to the new leaf. If the parent overflows too, it is split: the $j$-th value
$(j = \lfloor (p+1)/2 \rfloor)$ is \emph{moved} up to the grandparent, the first $j-1$ entries stay,
and the remaining entries go to a new internal node. The splitting may propagate up to the root; if
the root splits, a new root is created and the tree grows by one level. A B$^+$-tree therefore grows
from the top and stays balanced.

\paragraph{Example of insertion.}
Take $p = 3$ and $p_{leaf} = 2$ and insert the keys 8, 5, 1, 7, 3, 12, 9, 6 into an empty tree.
\Cref{tab:insert-trace} traces the evolution and \Cref{fig:bplus-final} shows the final tree.

\begin{table}[H]
\centering
\small
\begin{tabularx}{\textwidth}{@{}rX@{}}
\toprule
\textbf{Insert} & \textbf{Effect} \\
\midrule
8, 5 & Single leaf $[5,8]$ (the root is the leaf). \\
1 & Leaf overflows $[1,5,8]$: split into $[1,5]$ and $[8]$; $5$ is copied up, creating the root $[5]$. \\
7 & Goes to the right leaf: $[7,8]$. \\
3 & Left leaf overflows $[1,3,5]$: split into $[1,3]$ and $[5]$; $3$ is copied up, root becomes $[3,5]$. \\
12 & Right leaf overflows $[7,8,12]$: split into $[7,8]$ and $[12]$; $8$ is copied up and the root overflows $[3,5,8]$: $5$ moves up as the new root, $[3]$ and $[8]$ become internal nodes (tree height increases). \\
9 & Goes to the leaf $[12]$, giving $[9,12]$. \\
6 & Leaf $[7,8]$ overflows $[6,7,8]$: split into $[6,7]$ and $[8]$; $7$ is copied up, the internal node becomes $[7,8]$. \\
\bottomrule
\end{tabularx}
\caption{Insertion of 8, 5, 1, 7, 3, 12, 9, 6 with $p = 3$ and $p_{leaf} = 2$.}
\label{tab:insert-trace}
\end{table}

\begin{figure}[H]
\centering
\begin{tikzpicture}[
  nd/.style={draw, rounded corners=1pt, minimum height=7mm, inner sep=5pt, font=\small\ttfamily},
  >=Stealth, level distance=16mm]
  \node[nd, fill=blue!8] (r) at (0,0) {5};
  \node[nd, fill=blue!8] (a) at (-3.2,-1.6) {3};
  \node[nd, fill=blue!8] (b) at (3.2,-1.6) {7 \quad 8};
  \node[nd, fill=green!8] (l1) at (-4.6,-3.4) {1 \quad 3};
  \node[nd, fill=green!8] (l2) at (-1.8,-3.4) {5};
  \node[nd, fill=green!8] (l3) at (1.4,-3.4) {6 \quad 7};
  \node[nd, fill=green!8] (l4) at (3.4,-3.4) {8};
  \node[nd, fill=green!8] (l5) at (5.4,-3.4) {9 \quad 12};
  \draw[->] (r.south) -- (a.north);
  \draw[->] (r.south) -- (b.north);
  \draw[->] ([xshift=-2mm]a.south) -- (l1.north);
  \draw[->] ([xshift=2mm]a.south) -- (l2.north);
  \draw[->] ([xshift=-4mm]b.south) -- (l3.north);
  \draw[->] (b.south) -- (l4.north);
  \draw[->] ([xshift=4mm]b.south) -- (l5.north);
  \foreach \x/\y in {l1/l2,l2/l3,l3/l4,l4/l5} \draw[->, dashed] (\x.east) -- (\y.west);
\end{tikzpicture}
\caption{The B$^+$-tree after inserting 8, 5, 1, 7, 3, 12, 9, 6 ($p=3$, $p_{leaf}=2$); the leaves are linked in key order.}
\label{fig:bplus-final}
\end{figure}

\paragraph{Deletion.}
The entry is removed from the leaf; if the value also occurs in an internal node, the value to its
left in the leaf replaces it there. If the leaf then has fewer entries than the minimum
(\emph{underflow}), the DBMS first tries to \emph{redistribute} entries with a sibling leaf
(commonly the left sibling first, then the right). If redistribution is impossible, the leaf is
\emph{merged} with a sibling and the entry that pointed to the pair is deleted from the parent. An
internal node that underflows is handled in the same way, using the separating entry of the parent;
if the root becomes empty, the merged node becomes the new root. Merges can propagate up to the
root and reduce the number of levels.

\paragraph{Search.}
A point query follows one root-to-leaf path comparing the key with the node values, so it reads
as many blocks as there are levels plus one data block. A range query finds the first qualifying
leaf this way and then follows the leaf links, without returning to the upper levels.

\subsubsection{Index types in a commercial DBMS (Oracle, from class)}
The course also surveys the index types offered by Oracle, which serve as an example of
product technology.
\begin{itemize}
  \item \textbf{B-tree indexes} are the standard type; variants include index-organized tables
        (the data is itself the index), \emph{reverse-key} indexes (the bytes of the key are
        reversed to spread inserts over many blocks), \emph{descending} indexes and B-tree
        \emph{cluster} indexes.
  \item \textbf{Bitmap and bitmap join indexes} use a bitmap per value to point to many rows, and
        combine conditions with fast bit operations.
  \item \textbf{Function-based indexes} index a transformed column or an expression, for example
        \code{UPPER(name)}.
  \item \textbf{Application-domain indexes} are customized for a specific application.
\end{itemize}
Indexes are created with \code{CREATE [UNIQUE|BITMAP] INDEX \dots\ ON table(columns)}.

\subsubsection{Key points used in the comparison}
\begin{enumerate}
  \item An index is an ordered access path of $\langle K, P\rangle$ entries: it speeds up reads at
        the price of extra storage and slower inserts, deletes and updates.
  \item Primary and clustering indexes are sparse and require the data file to be physically
        ordered (hence only one of them per file), whereas secondary indexes are dense and can be
        created on any field, in any number.
  \item Multilevel indexes reduce the search cost to $t+1$ block accesses, where $t$ is logarithmic
        in the number of entries with the fan-out as the base.
  \item B$^+$-trees keep the index balanced under insertions and deletions by splitting and merging
        nodes, store data pointers in linked leaves, and are the structure used in practice.
\end{enumerate}

\subsection{Indexing in MongoDB (practice)}
\label{sec:index-practice}
\todo{WiredTiger B-tree, \code{\_id} index, key $\rightarrow$ RecordId, clustered collections, compound / multikey / hashed / vector indexes.}

\subsection{Demo and results}
\label{sec:index-demo}

\subsubsection{Case 1: retrieve by patient ID}
\todo{COLLSCAN vs IXSCAN, explain() metrics.}

\begin{lstlisting}[language=mongosh,caption={Example: create an index and inspect the plan}]
db.images.createIndex({ patient_id: 1 })
db.images.find({ patient_id: "P001" }).explain("executionStats")
\end{lstlisting}

\subsubsection{Case 2: patient ID with conditions}
\todo{Compound index, ESR rule, index prefix, covered query.}

\subsubsection{Case 3: similar MRI images}
\todo{Embeddings, vector search index, ANN vs exact.}

\subsection{Theory vs. practice}
\label{sec:index-compare}
\begin{table}[H]
\centering
\begin{tabularx}{\textwidth}{@{}>{\raggedright\arraybackslash}X>{\raggedright\arraybackslash}X@{}}
\toprule
\textbf{Class (Chapter 2)} & \textbf{MongoDB} \\
\midrule
\todo{fill in} & \todo{fill in} \\
\bottomrule
\end{tabularx}
\caption{Comparison between theory and practice.}
\end{table}
